\section{Implementation}

Now that we know what the IR is, we shall present in this section the
structure and infrastructure provided by \MLIR and how it is used in practice.

% ODS, passes (constant prop, CSE, inlining, etc can be just listed), Interfaces  with inlining as an example, convey the idea of multiple dialects mixed together and why that is a good thing.

\subsection{Structure}

\begin{figure}
	% TODO: use example from tutorial slides.
\lstinputlisting[basicstyle=\tiny]{example.mlir}
\caption{\MLIR generic textual format for Op.}
\end{figure}

\MLIR is a compiler intermediate representation (IR) based on static single
assignment (SSA) form.  It takes a radical approach to the simplicity and
extensibility of compiler IRs. The unit of semantics in \MLIR is an operation,
commonly referred to as an Op. Ops are are equivalent to instructions in
conventional IRs, but can also provide additional structure as described below.
\MLIR does not have a predefined set of Ops.  Instead, it operates on sets of
Ops with user-defined semantics. This means that users of \MLIR can mix and
match existing Ops and introduce new ones to suit their needs. The core
infrastructure is designed to conservatively operate on opaque Ops whose
semantics may be unknown. In particular, opaque Ops can be parsed, printed and
copied without any knowledge of their semantics.

\subsubsection{Operations}

Ops accept as operands and produce as results lists of typed values, which obey
SSA. Values represent data that the IR operates upon, i.e. the data that is
known at run time. In addition to that, an Op may have a set of attributes with
values known at compile time. On the structure side, Ops may have successor
blocks as well as attached regions.

\subsubsection{Blocks and CFG}

Lists of Ops form (basic) blocks. The last Op in a block is always a terminator
Op, which may or may not have successors. Terminators are commonly used to
describe control across blocks and are, together with regions, the core control
flow abstractions. The transfer of control is up to the semantics of the
specific terminator Op. For example, a terminator may transfer control to
another named block (branch operation), or immediately exit from a surrounding
region (return operation).

References to successor blocks in a terminator Op define a control flow graph
(CFG), where nodes correspond to blocks and edges correspond to the links
between a block and successors listed in its terminator Op. \MLIR adopts block
arguments as the way to pass values across blocks, in contrast to the often
used phi instructions. Each block has a potentially empty list of typed block
arguments, which are regular values and obey SSA. Terminator Ops referring to
successor blocks are required to pass values to block arguments, which is
roughly equivalent to renaming the value inside before using it in the next
block. For the purposes of SSA, block arguments are considered to be defined
before the first Op in the block.

Blocks are a core IR concept and provide one level of value scoping. Ops in a
block can be thought of as executed in order. As a result, any operation is
allowed to use as operands the values produced by preceding operations, in
addition to block arguments and other visibility rules based on dominance and
structure (see below).

\subsection{Regions}
Lists of blocks form regions, which are an essential structuring abstraction in
\MLIR. Regions do not have semantics by themselves. It is assigned to them by
the Op they are associated with. The containing Op defines if and how a region
is executed.

This approach enables straightforward definition of structured control flow
using regions. A conditional (\lstinline!if! operation) is merely an Op with a
condition operand and two regions: one containing the ``true'' branch and another
containing the ``false'' branch. The semantics of the op could define that the
control flow is transferred to the first or the second region based on the
value of the operand.

The first block of a region is referred to as entry block and can be
semantically connected to the surrounding Op. In particular, the Op semantics
prescribes what values the arguments of the entry block have when the control
flow is transferred to the region. The arguments of the entry block are treated
as arguments of the region containing the block. These values can be connected
to the Op operands. For example, a simple counted loop (\lstinline!for! operation) can be
implemented as an Op with a single region executed multiple times with an
argument taking the value of the induction variable on each step.

\subsubsection{Symbols and Symbol Tables}
<...>

\subsubsection{Functions and Modules}
Similarly to conventional IRs, \MLIR can be structured into functions and
modules. However, these are not first-class concepts in \MLIR. Instead they are
implemented as Ops.

A module is an Op with a single region containing a single block, and
terminated by a dummy Op that does not transfer the control flow. A module is a
symbol and can be referenced. Like any block, its body contains a list of Ops,
which may or may not have the semantics of a function. For example, modules may
also contain Ops for global variables or compiler metadata.

A function is an Op with a single region, with arguments corresponding to
function arguments. It is a symbol and can be referenced. The control flow is
transferred into a function using a function call Op. Once inside, the control
flow follows the CFG of the blocks in the region. A special terminator returns
the control flow from the body region and transfers it to the immediate
control-flow successor of the call Op. Any operands of the ``return'' terminator
are transmitted to the result of the call Op.

\subsubsection{Type System}

\subsection{Operation description}

Ops form the core of \MLIR, and \MLIR allows and open Op ecosystem, but Ops
with regions and verification enables very powerful abstractions. In \MLIR we
wanted to make it easy to define operations and their verification, to this end
we developed Operation Descriptions (ODS) a declarative format to define Ops
using the data modelling language TableGen~\cite{tblgen}.

\begin{figure}
\lstinputlisting[basicstyle=\tiny]{leakyrelu.td}
\caption{Operation description of LeakyRelu.}
\end{figure}

\subsubsection{Attributes}

\subsubsection{Traits}

\subsubsection{Verifier}

Ops can have verifiers defined and these verifiers can also act upon the
regions of Ops. This allows verifying invariants from local properties, e.g.,
the input operands to the Op are integral, to more structural, e.g., the
accesses within the affine loop are indeed affine.

\subsection{Declarative rewrites}

Ops are the fundamentally representational element in \MLIR, and manipulating
Ops core to all transformations and lowerings in \MLIR. Some transformations
can be expressed simply as $N-M$ DAG rewrites while others require more
complicated transformations. In \MLIR we did not want to give up either the
generality of op transformations or the simplicity of DAG-DAG rewrites, so we
developed a general graph rewriting framework and layered a declarative rewrite
system on top of it. Thereby not complicating the simple case for the sake
of generality.

\subsection{Passes}

\subsection{Interfaces}\label{sec:interfaces}

Passes in \MLIR are intended to apply to multiple dialects and multiple levels,
interfaces provides the mechanism to enable this flexibility. Interfaces
provide a mechanism to define passes operating on structure of operations and
allowing dialects/Ops to opt-in to specific optimization behavior.

Consider as an example inlining.
